\section{Metric Foundations of Exact Information Reduction}
\label{sec:metric-foundations}

We first separate exact information reduction from smooth or matrix-specific
structure.  Let \((X,d_X)\) and \((Y,d_Y)\) be arbitrary metric spaces and let
\(q:X\to Y\) be surjective.

\begin{definition}[Metric submetry]
\label{def:metric-submetry}
The map \(q\) is a metric submetry if for every \(x\in X\) and \(r\ge0\),
\begin{equation}
  q(\overline B_X(x,r))=\overline B_Y(q(x),r).
  \label{eq:metric-submetry}
\end{equation}
\end{definition}

A submetry maps every closed ball onto the ball of the same radius.  It is
therefore stronger than a nonexpansive surjection: it guarantees that every
visible displacement has an ambient lift with no excess distance.

\begin{theorem}[Maximal exact information reduction]
\label{thm:metric-characterization}
The following statements are equivalent.
\begin{enumerate}[label=(\Alph*),leftmargin=*]
  \item \(q\) is a metric submetry.
  \item For every \(x\in X\) and \(y\in Y\), the fiber distance is attained
  and exact:
  \begin{equation}
    d_Y(q(x),y)=\min_{q(z)=y}d_X(x,z).
    \label{eq:exact-fiber-distance}
  \end{equation}
  \item For every nondecreasing
  \(\rho:[0,\infty)\to\R\cup\{+\infty\}\), every
  \(\phi:Y\to\R\cup\{+\infty\}\), and every \(x\in X\),
  \begin{equation}
    \boxed{
    \inf_{z\in X}\{\rho(d_X(x,z))+\phi(q(z))\}
    =
    \inf_{y\in Y}\{\rho(d_Y(q(x),y))+\phi(y)\}.}
    \label{eq:universal-radial-reduction}
  \end{equation}
\end{enumerate}
Thus metric submetries are exactly the observation maps for which every
reference-radial, visible-functional decision problem reduces without loss.
\end{theorem}

The proof, given in \cref{app:metric-proofs}, is elementary but sharp.  Hard
ball costs and singleton visible losses recover the full ball-surjectivity
condition, so the theorem cannot be enlarged while preserving universal exact
reduction.

The theorem determines optimal values but does not choose one lift when a
fiber contains several nearest points.  Canonical completion requires a
coherent family of sections.  Suppose \(q\) is \(1\)-Lipschitz and for every
\(x\in X\) a section \(s_x:Y\to X\) is given such that
\begin{equation}
  q\circ s_x=\operatorname{id}_Y,
  \qquad
  s_x(q(x))=x,
  \qquad
  s_{s_x(y)}=s_x.
  \label{eq:coherent-sections}
\end{equation}
The last identity says that moving inside one selected information sheet does
not change the sheet.

For a function \(F\) on a metric space, we use
\[
  e_\tau F(x)=\inf_z\left\{F(z)+\frac{d(x,z)^2}{2\tau}\right\}
\]
for its Moreau envelope and \(\prox_{\tau F}(x)\) for the corresponding
minimizer when it is unique.

\begin{theorem}[Coherent exact reduction]
\label{thm:coherent-reduction}
Under \eqref{eq:coherent-sections}, the following are equivalent:
\begin{enumerate}[label=(\Alph*),leftmargin=*]
  \item every \(s_x\) is an isometric embedding;
  \item for every \(x\in X\) and \(y\in Y\),
  \begin{equation}
    d_Y(q(x),y)
    =d_X(x,s_x(y))
    =\min_{q(z)=y}d_X(x,z);
    \label{eq:coherent-fiber-distance}
  \end{equation}
  \item the equality \eqref{eq:universal-radial-reduction} holds and
  \(s_x\) maps every reduced minimizer to a full minimizer.
\end{enumerate}
Whenever these conditions hold, \(q\) is a metric submetry.  If the nearest
point in every fiber is unique, the splitting is rigid.

If \(Y\) is complete CAT(0), \(\phi\) is proper, lower semicontinuous,
geodesically convex, and bounded below, and \(\tau>0\), then
\begin{equation}
  \boxed{
  \prox^{X,\mathrm{can}}_{\tau(\phi\circ q)}(x)
  =
  s_x\!\left(\prox^Y_{\tau\phi}(q(x))\right).}
  \label{eq:coherent-prox}
\end{equation}
Coherent proximal iterations are therefore exact lifts of the visible
iterations; under rigidity, the full proximal point is unique.
\end{theorem}

\begin{remark}[Three distinct requirements]
Surjectivity provides visible representatives.  Submetry provides exact
distance.  Coherence provides a reusable canonical sheet.  None of these
properties follows from the others without the stated assumptions.
\end{remark}
